\documentclass[12pt]{article}
\usepackage[utf8]{inputenc}
\usepackage[T1]{fontenc}
\usepackage{hyperref}
\title{Sasaru Language Resources for K-12 Curriculum Development}
\author{Compiled by Claude Sonnet 5.5 (claude-sonnet-5-5) for Andr\'as Kornai}
\date{}
\begin{document}
\maketitle

\section{Introduction}
Sasaru (ISO 639-3 sxs) is an Edoid language of Akoko-Edo, Edo State, Nigeria. Speaker figures conflict: Wikipedia cites 13,000 (2000), and Joshua Project lists a people-group population of 26,000. Glottolog groups Sasaru with Igwe. Very little material was found, and nothing educational.

\section{Primary Educational Resources}
No primers, textbooks or school materials for Sasaru were found.

\section{Linguistic and Reference Materials}
\begin{itemize}
\item Wikipedia, ``Sasaru language'': \url{https://en.wikipedia.org/wiki/Sasaru_language}
\item Glottolog language map for Sasaru: \url{https://glottolog.org/resource/languoid/id/sasa1252.bigmap.html}
\item Glottolog language map for Igwe: \url{https://glottolog.org/resource/languoid/id/igwe1238.bigmap.html}
\end{itemize}
No dictionary or grammar was found.

\section{Religious and Community Resources}
\begin{itemize}
\item Joshua Project, Sasaru people profile (Scripture status: translation started; no audio Bible or Jesus Film listed): \url{https://joshuaproject.net/people_groups/14777}
\item Joshua Project, language page for sxs: \url{https://joshuaproject.net/languages/sxs}
\end{itemize}

\section{Digital and Audio Resources}
No audio, video or software resources were found.

\section{Curriculum Development Context}
The language is not known to be taught. A curriculum would start from basic documentation (word lists, orthography, recordings) and would need to be built with the Sasaru community in Akoko-Edo. The sources above are reference-level only.

\end{document}